\section{Production and accounts (to assign)}\label{sec:accounts}

Production and accounts turn sales into output, costs, profits and balance sheets; every
other mechanism reads them, and this is where the model checks its stock-flow consistency.

\paragraph{The trap: the accounts define the profit rate every rule compares.} Profit gross
of depreciation puts every rate $\delta$ above the net rate $(p - uc)/\kappa$ that $\rbar$
measures; dividing by book value makes the oldest plants look the best. One definition is
used everywhere: profit net of depreciation and interest, over capital at purchase cost.
Book values serve only for net worth.

\paragraph{Starting rule.} Firms produce to order: output equals sales, and plants run in
merit order, lowest operating cost first. Output above normal capacity costs a premium of
10\,\% on $ucs$ (\code{OVERTIME\_PREMIUM}), the same for both types.
\begin{align*}
\Pi_{i,t} &= (p_{i,t} - ucs_i)\,x_{i,t} - \text{premium} - \textstyle\sum_v \delta\,\kappa_v\,\qbar_v - \rho\,B_{i,t-1},
  \qquad r_{i,t} = \frac{\Pi_{i,t}}{\sum_v \kappa_v\,\qbar_v}, \\
M_{i,t} &= M_{i,t-1} + \Pi_{i,t} + \textstyle\sum_v \delta\,\kappa_v\,\qbar_v - \mathit{Div}_{i,t}
  - \kappa_{j^*}\,I_{i,t} + \Delta B_{i,t}, \qquad \mathit{Div}_{i,t} = d\,\max\{0, \Pi_{i,t}\}, \\
NW_{i,t} &= M_{i,t} + K^{b}_{i,t} - B_{i,t},
\end{align*}
with $\rho = 0.05 < \rbar$ and $d = 0.5$ for both types. Book values fall in a straight line
over $T$ periods; a plant retired early leaves its book value as a write-off, which lowers
net worth while no cash moves. Firms repay 10\,\% of their debt a period from cash. In version 1 a
negative cash balance is covered by an overdraft loan; in version 2
(\code{CASH\_SHORTFALL = "bankruptcy"}, the default) it stays as bills the firm could not
pay, and the firm goes bankrupt at the end of the period (section~\ref{sec:entry}). Since a
firm sells only what plants whose $ucs$ the price covers can make, and dividends come out of
positive profit, cash runs out only when interest exceeds the operating margin.

\paragraph{Alternative: depreciation by use.} With \code{DEPRECIATION = "use"} (the default
is \code{"time"}) a plant depreciates with what it produces, $\delta\,\kappa_v\,x_{v,t}$
(overtime included), instead of $\delta\,\kappa_v\,\qbar_v$ a period, until its book value is
zero. Its purchase cost is then spread over a lifetime output of $\qbar_v/\delta$. At normal
output both rules charge the same, and the capital charge per unit of output is
$\delta\kappa_v$ in both, so $uc = ucs + \delta\kappa$, and with it the hurdle, the
technique choice, replacement and the scrap cost, are unchanged. What changes is the cost of
idle capacity: an idle plant does not depreciate, so a firm at low utilisation reports a
higher profit, and a plant scrapped while idle leaves a larger write-off. Under this rule
\code{PLANT\_LIFE} counts periods of normal output rather than age: a plant leaves once it
has produced \code{PLANT\_LIFE}$\times\qbar_v$, which with \code{PLANT\_LIFE} $= 1/\delta$ is
when it is fully depreciated. The rule is the same for both types.

\paragraph{Closure.} Wages and materials, plant purchases, interest, dividends, loans,
equity and the industry's demand all go through one rest-of-economy account. Every payment
goes through \code{markets/payments.pay}, which records the transaction-flow matrix. At the
end of each period the scheduler asserts that money is conserved, that firms' debt equals
the loans of the rest of the economy, that each firm's net worth moved by
$\Pi - \mathit{Div} - \text{write-offs}$ (plus equity at entry), and that every row and
column of the flow matrix adds up.

\paragraph{Tests} (\code{tests/test\_accounts.py}, \code{tests/test\_production.py}): at
normal output with no debt $r_i = (p - uc)/\kappa$, i.e.\ 16.00\,\%, 16.86\,\% and 15.83\,\%
at $p = 100$; the net-worth identity; the flow matrix sums to zero; depreciation charges add
up to the purchase cost and book value is zero at age $T$; a firm without sales loses
exactly depreciation plus interest; merit order, overtime and ties; label swap.

\subsection*{Code}
\eqfile{merit_order.py}
\eqfile{operating_cost.py}
\eqfile{depreciation.py}
\eqfile{interest_payment.py}
\eqfile{firm_profit.py}
\eqfile{profit_rate.py}
\eqfile{plant_profit.py}
\eqfile{dividends.py}
\eqfile{debt_repayment.py}
\marketfile{payments.py}
\marketfile{credit_market.py}
\schedulerpart{accounts}{\code{model/scheduler.py}: step 3}
